\documentclass[11pt]{article}
\usepackage[margin=1in]{geometry}
\usepackage{amsmath,amssymb}
\usepackage{booktabs}
\usepackage{hyperref}
\title{Detailed Technical Notes for a Darts-Based LTSF Benchmark}
\date{}
\begin{document}
\maketitle

\section{Benchmark scope and relation to arXiv:2205.13504v3}
This benchmark follows the same nine datasets used in arXiv:2205.13504v3: ETTh1, ETTh2, ETTm1, ETTm2, Electricity, Traffic, Weather, Exchange-Rate, and ILI. For each dataset, the benchmark evaluates AutoARIMA, Linear, NLinear, DLinear, LSTM, GRU, N-BEATS, and Transformer using Darts, and reports MAE, MSE, training time, and inference time.

Let an input window be $\mathbf{x}_{t-L+1:t}\in\mathbb{R}^{L\times C}$ with look-back length $L$ and $C$ channels, and the forecast horizon be $H$. Each model produces
\[
\hat{\mathbf{y}}_{t+1:t+H} = f_\theta(\mathbf{x}_{t-L+1:t}), \quad \hat{\mathbf{y}}\in\mathbb{R}^{H\times C}.
\]

\section{AutoARIMA}
\subsection{Core model}
ARIMA$(p,d,q)$ models one scalar time series (or one channel) after differencing order $d$:
\[
\phi(B)(1-B)^d y_t = c + \theta(B)\varepsilon_t,
\]
where $B$ is the backshift operator, $\phi(B)=1-\phi_1B-\cdots-\phi_pB^p$, and $\theta(B)=1+\theta_1B+\cdots+\theta_qB^q$.

\subsection{Why ``Auto''}
AutoARIMA performs order selection over candidate $(p,d,q)$ (and seasonal terms if enabled), typically minimizing AIC/AICc/BIC while checking stationarity/invertibility constraints. The benchmark uses it as a classical non-neural baseline. For multivariate series, a separate model is fit per channel and channel forecasts are stacked.

\subsection{Strengths and limitations}
AutoARIMA is data-efficient and interpretable for linear autocorrelation structures, but can underfit strong nonlinearities and long multivariate cross-channel interactions.

\section{Linear (lagged linear regression)}
The Linear baseline uses lagged values as features and predicts a multi-step target with a linear map:
\[
\mathrm{vec}(\hat{\mathbf{y}})=\mathbf{W}\,\mathrm{vec}(\mathbf{x})+\mathbf{b}.
\]
In Darts this is implemented as a regression model with lag features. Training is convex (least-squares style objective under default settings), giving stable optimization and fast fitting/inference.

This baseline answers an important question: how far can purely linear temporal projection go before architectural complexity is needed?

\section{NLinear}
NLinear (Normalization-Linear) reduces local level shift before linear mapping. For each channel, define the last observed value in the look-back window as $x_t$. Normalize:
\[
\tilde{\mathbf{x}} = \mathbf{x} - x_t\mathbf{1}_L.
\]
Apply a learned linear projection $g_\theta$ to $\tilde{\mathbf{x}}$, then restore level:
\[
\hat{\mathbf{y}} = g_\theta(\tilde{\mathbf{x}}) + x_t\mathbf{1}_H.
\]

This explicitly targets distribution shift in absolute level between train and test windows. Compared with plain Linear, NLinear changes the inductive bias from ``fit absolute scale'' to ``fit local deviations from current level.''

\section{DLinear}
DLinear (Decomposition-Linear) first decomposes input into trend and seasonal/remainder components using moving average:
\[
\mathbf{x}^{\text{trend}}=\mathrm{MA}(\mathbf{x}),\qquad
\mathbf{x}^{\text{season}}=\mathbf{x}-\mathbf{x}^{\text{trend}}.
\]
Then two separate linear heads are learned:
\[
\hat{\mathbf{y}}=g^{\text{trend}}_\theta(\mathbf{x}^{\text{trend}})+g^{\text{season}}_\theta(\mathbf{x}^{\text{season}}).
\]

The decomposition imposes a useful bias: low-frequency trend and high-frequency residual dynamics are modeled with different parameter sets. This often improves robustness on data where trend drift and periodic components coexist.

\section{RNN family: LSTM and GRU}
RNNs model temporal state recursively:
\[
\mathbf{h}_t=\Psi(\mathbf{x}_t,\mathbf{h}_{t-1}).
\]
They share parameters across time and naturally support variable sequence lengths.

\subsection{LSTM}
LSTM introduces memory cell $\mathbf{c}_t$ and gates:
\begin{align*}
\mathbf{i}_t &= \sigma(\mathbf{W}_i[\mathbf{x}_t,\mathbf{h}_{t-1}] + \mathbf{b}_i),\\
\mathbf{f}_t &= \sigma(\mathbf{W}_f[\mathbf{x}_t,\mathbf{h}_{t-1}] + \mathbf{b}_f),\\
\mathbf{o}_t &= \sigma(\mathbf{W}_o[\mathbf{x}_t,\mathbf{h}_{t-1}] + \mathbf{b}_o),\\
\tilde{\mathbf{c}}_t &= \tanh(\mathbf{W}_c[\mathbf{x}_t,\mathbf{h}_{t-1}] + \mathbf{b}_c),\\
\mathbf{c}_t &= \mathbf{f}_t\odot\mathbf{c}_{t-1} + \mathbf{i}_t\odot\tilde{\mathbf{c}}_t,\\
\mathbf{h}_t &= \mathbf{o}_t\odot\tanh(\mathbf{c}_t).
\end{align*}
The additive memory update helps gradient flow across long contexts.

\subsection{GRU}
GRU simplifies gating with no separate cell state:
\begin{align*}
\mathbf{z}_t &= \sigma(\mathbf{W}_z[\mathbf{x}_t,\mathbf{h}_{t-1}] + \mathbf{b}_z),\\
\mathbf{r}_t &= \sigma(\mathbf{W}_r[\mathbf{x}_t,\mathbf{h}_{t-1}] + \mathbf{b}_r),\\
\tilde{\mathbf{h}}_t &= \tanh(\mathbf{W}_h[\mathbf{x}_t,\mathbf{r}_t\odot\mathbf{h}_{t-1}] + \mathbf{b}_h),\\
\mathbf{h}_t &= (1-\mathbf{z}_t)\odot\mathbf{h}_{t-1}+\mathbf{z}_t\odot\tilde{\mathbf{h}}_t.
\end{align*}
GRU often trains faster than LSTM for similar hidden widths.

\subsection{Forecasting behavior}
In Darts RNN forecasting, models are trained on sliding windows and produce multi-step outputs autoregressively or chunk-wise depending on configuration. They are strong for smooth temporal dynamics but can be compute-heavy for very long horizons.

\section{N-BEATS}
N-BEATS uses stacks of fully-connected residual blocks. Each block receives an input representation $\mathbf{z}^{(k)}$ and outputs:
\[
\mathbf{b}^{(k)} \; (\text{backcast}),\qquad
\mathbf{f}^{(k)} \; (\text{forecast}).
\]
Residual recursion:
\[
\mathbf{z}^{(k+1)} = \mathbf{z}^{(k)} - \mathbf{b}^{(k)},
\]
and total forecast:
\[
\hat{\mathbf{y}} = \sum_k \mathbf{f}^{(k)}.
\]

Intuition: each block explains part of the input history (backcast) and contributes forecast signal; later blocks model residual structure left by earlier blocks. This often gives strong accuracy without recurrent or attention modules.

\section{Transformer}
\subsection{Attention mechanism}
Given token embeddings $\mathbf{X}$, self-attention forms queries/keys/values:
\[
\mathbf{Q}=\mathbf{X}\mathbf{W}_Q,\quad
\mathbf{K}=\mathbf{X}\mathbf{W}_K,\quad
\mathbf{V}=\mathbf{X}\mathbf{W}_V,
\]
then computes
\[
\mathrm{Attn}(\mathbf{Q},\mathbf{K},\mathbf{V})=\mathrm{softmax}\!\left(\frac{\mathbf{Q}\mathbf{K}^\top}{\sqrt{d_k}}\right)\mathbf{V}.
\]
Multi-head attention concatenates several such projections, followed by feed-forward sublayers, residual connections, and normalization.

\subsection{Forecasting interpretation}
Transformer forecasting maps input windows to horizon outputs while using positional information to preserve temporal order. It is expressive for long-range interactions, but has higher optimization and compute cost than linear baselines on many LTSF settings.

\section{Metrics and timing protocol}
For each dataset--horizon--model run:
\begin{itemize}
    \item \textbf{MAE}: $\frac{1}{N}\sum_i |y_i-\hat y_i|$.
    \item \textbf{MSE}: $\frac{1}{N}\sum_i (y_i-\hat y_i)^2$.
    \item \textbf{train\_time}: wall-clock time spent in model fitting.
    \item \textbf{inference\_time}: wall-clock time for one $H$-step prediction pass.
\end{itemize}
All run-level rows are written to CSV and aggregated summaries are computed over successful runs only.

\end{document}
